\fichatitulo{05 · SÍNTESE TEMÁTICA}{Interpretação do piloto}{O que as leituras articulam}
Esta síntese relaciona as três fichas por tema. As propostas para a dissertação são interpretações de trabalho, sujeitas à análise do pesquisador. Os códigos [A], [K] e [I] remetem às fontes e aos localizadores da matriz.

\campo{1. Qualidade exige definir o objeto da avaliação}
As fontes permitem distinguir a avaliação do componente, a apresentação da aplicação e a formulação de expectativas institucionais. Esses objetos se complementam, mas não são intercambiáveis: uma descrição funcional não estabelece uma medida de desempenho, e uma medida técnica não responde sozinha à pergunta sobre utilidade institucional ([A], seções 3--5; [K], seções 2--3; [I], Introdução e caso 024).

\textbf{Implicação proposta:} explicitar na metodologia o que cada instrumento observa. A conexão entre relevância, fidelidade e verificabilidade precisa ser definida para o nosso caso, sem transformar os três termos em sinônimos.

\campo{2. A participação humana assume funções diferentes}
Anotação de referência, retorno de usuários e satisfação cumprem papéis distintos na comparação. A aproximação temática não autoriza reuni-los em uma única categoria de validação humana ([A], seções 3.3 e 5.1; [K], seção 3; [I], caso 024, Success metrics).

\textbf{Implicação proposta:} separar, no protocolo, julgamento de qualidade documental e experiência do usuário. O primeiro se relaciona diretamente com H3; o segundo demandaria um instrumento próprio caso venha a integrar a pesquisa.

\campo{3. Há complementaridade, sem contradição demonstrada}
Neste recorte, as fontes não oferecem testes equivalentes que permitam decidir qual sistema é melhor. As diferenças decorrem principalmente de finalidade, desenho e evidência apresentada. Não foi identificada uma contradição empírica comparável; isso não significa que a literatura mais ampla seja consensual ([A], seção 4; [K], seções 2--4; [I], Introdução).

\textbf{Implicação proposta:} usar os textos para justificar critérios de comparação e perguntas a investigar. O fato de uma ficha não conter determinado resultado não basta para afirmar originalidade ou lacuna de pesquisa.

\campo{Questões para orientar a continuidade da revisão}
\begin{itemize}
\item Quais estudos avaliam conjuntamente recuperação, geração e suporte documental de afirmações em acervos parlamentares?
\item Como são construídos os conjuntos de referência e tratadas as divergências entre julgamentos humanos e automatizados?
\item Há evidência comparável em português e em condições próximas às do corpus da dissertação?
\end{itemize}

\campo{Próxima etapa bibliográfica}
Aplicar a matriz às demais referências selecionadas e documentar a busca complementar. Cada afirmação de convergência, divergência ou lacuna deverá remeter às passagens pertinentes e respeitar o alcance dos desenhos estudados. A quantidade de referências concordantes, isoladamente, não mede a força da evidência.
